\chapter{Dividendes, emprunt et forwards}\label{ch:dv:dividends-borrow-and-forwards}

Une action valant 50 figure sur la liste des titres difficiles à emprunter de tous les courtiers. Ses options à un mois semblent faussées~:
à la valeur du strike de 50, le put s'échange à 4.04 et le call à 2.78,
et un écran qui calcule les volatilités implicites à partir du spot et du taux d'intérêt
affiche une volatilité de 72~\% pour le put et de 47~\% pour le call. Un nouveau trader y voit 25
points d'arbitrage~: vendre le put cher, acheter le call bon marché, couvrir. Il n'y en a
aucun. La parité call-put (livre~1, chapitre~25) dit que l'écart entre le call et le put est un forward,
et le forward que les options intègrent est
48.74, et non les 50.16 que suggèrent le spot et le taux~: le marché facture
35~\% par an pour emprunter le titre, et quiconque veut être vendeur à découvert
par le biais des options paie ce taux dans les deux sens. Ce chapitre porte sur le forward~: ce qui
entre dans son prix pour une action (les dividendes et le coût de l'emprunt), comment les dividendes
entrent dans un modèle d'option, comment la chaîne d'options elle-même révèle le forward, et qui porte
le risque que les dividendes changent.

\section{Forwards avec dividendes discrets et emprunt}

Une action verse des dividendes à des dates de détachement connues (livre~1, chapitre~8),
et un vendeur à découvert verse des frais d'emprunt au prêteur (livre~1, chapitre~6).
Ces deux éléments font baisser le forward sous le spot financé~: l'acheteur d'un
forward ne touche pas les dividendes versés avant la livraison, et le vendeur qui détient l'action
comme couverture peut la prêter et en encaisser les frais.

\begin{proposition}[Le forward sur action]\label{prop:dv:dividends-borrow-and-forwards:fwd}
Avec des dividendes en espèces $D_i$ versés en $t_i$, un taux de financement plat $r$ et des frais d'emprunt
$\ell$ (capitalisés en continu, payés sur la valeur prêtée), le forward de livraison
en $T$ vaut
\[
F_{0,T}=\Bigl(S_0-\sum_{t_i\le T}D_iP(0,t_i)\Bigr)\frac{e^{-\ell T}}{P(0,T)} .
\]
\end{proposition}
\begin{proof}
Acheter l'action au prix $S_0$, emprunter la valeur actuelle des dividendes pour les financer,
prêter l'action et réinvestir les frais, vendre le forward. En $T$, la position est
neutre et les dividendes ont remboursé leur emprunt~: le coût de l'action, diminué de ce
qu'elle rapporte, doit égaler la valeur actuelle du prix de livraison.
\end{proof}

\begin{definition}[Dividende proportionnel]\label{def:dv:dividends-borrow-and-forwards:prop}
Un \emph{dividende proportionnel}\index{dividende proportionnel} se modélise comme une fraction fixe
$y_j$ du cours de l'action à la date de détachement, et non un montant fixe~: l'action
recule de $y_jS_{t_j-}$, et le forward est multiplié par $(1-y_j)$ pour chaque
date de détachement antérieure à $T$. Un rendement de dividende continu $q$ en est la limite lorsque les versements
sont répartis régulièrement dans le temps.
\end{definition}

La courbe à terme d'une action est une dent de scie
(\cref{fig:dv:dividends-borrow-and-forwards:curve})~: elle croît au taux de financement
diminué de l'emprunt entre les dates de détachement, et chute de chaque dividende. Un rendement continu
avec le même forward à deux ans est une droite qui passe par ces points, fausse à chaque
date intermédiaire d'un montant pouvant atteindre la moitié du dividende trimestriel. Pour les options dont l'échéance
chevauche une date de détachement, l'écart est le dividende entier.

\begin{omfigure}
\begin{tikzpicture}
\begin{axis}[width=10.5cm,height=4.8cm,
  xlabel={date de livraison (années)},ylabel={forward},
  tick label style={font=\small},label style={font=\small},
  xmin=0,xmax=2,ymin=99.3,ymax=100.9,grid=major,grid style={gray!25},
  yticklabel style={/pgf/number format/.cd,fixed,precision=1},
  legend columns=2,legend style={font=\footnotesize,at={(0.5,-0.3)},anchor=north,draw=none,fill=none,/tikz/every even column/.append style={column sep=8pt}},
  table/col sep=comma]
\addplot[omDef,very thick,const plot mark left] table[x=t,y=discrete]{figdata/derivatives/05-dividends-borrow-and-forwards/forward_curve.csv};
\addplot[omThm,thick,dashed] table[x=t,y=smooth]{figdata/derivatives/05-dividends-borrow-and-forwards/forward_curve.csv};
\legend{dividendes trimestriels en espèces de 0.60,rendement continu avec le même forward à 2 ans}
\end{axis}
\end{tikzpicture}

\omcaption{Courbe à terme d'une action valant 100 et versant 0.60 par trimestre, avec un taux
de financement de 3~\% et des frais d'emprunt de 0.5~\%. Entre les dates de détachement, le
forward croît de 2.5~\% par an~; à chaque date de détachement, il chute du montant du dividende. Paramètres
illustratifs.}\label{fig:dv:dividends-borrow-and-forwards:curve}
\end{omfigure}

\section{Dividendes dans le modèle d'option}

Pour une option européenne, seuls comptent le forward et la distribution qui l'entoure~; la
formule de Black du chapitre~3 la valorise sur le forward. Le choix du modèle de dividende
détermine de quelle volatilité la volatilité cotée est la volatilité, et comment le forward
évolue lorsque le spot bouge.

\begin{definition}[Modèle à dividendes séquestrés]\label{def:dv:dividends-borrow-and-forwards:escrowed}
Dans le \emph{modèle à dividendes séquestrés}\index{modèle à dividendes séquestrés}, l'action se
décompose en la valeur actuelle des dividendes à verser avant l'échéance, traités comme
sans risque, et le reste, $S^*_t=S_t-\sum_{t<t_i\le T}D_iP(t,t_i)$, qui suit un
mouvement brownien géométrique. Les options européennes se valorisent par Black--Scholes avec
$S^*_0$ à la place de $S_0$.
\end{definition}

Trois modèles sont utilisés. Dans le modèle séquestré, la volatilité est celle de $S^*$.
Dans le modèle proportionnel, elle est celle de l'action entière, qui reste lognormale
à travers ses baisses proportionnelles. Dans le \emph{modèle spot}, l'action entière est
lognormale entre les dates de détachement et chute du montant en espèces à chaque date~: le plus
réaliste, et le seul sans formule fermée. Pour une option européenne à même forward,
les modèles séquestré et proportionnel donnent le même prix à même
volatilité, puisqu'ils rendent tous deux le forward lognormal avec cette volatilité~; le modèle spot
donne un prix plus élevé, car avant la date de détachement, la part de l'action qui
sera versée en dividende bouge elle aussi.

\begin{proposition}[La volatilité de l'action séquestrée]\label{prop:dv:dividends-borrow-and-forwards:adj}
Si l'action a une volatilité $\sigma$ dans le modèle spot et verse un dividende de
valeur actuelle $\bar D$ en $t_1<T$, le modèle séquestré reproduit son prix à la monnaie
au premier ordre avec la volatilité effective
\[
\sigma_{\mathrm{eff}}^2\,T=\sigma^2\Bigl(\frac{S_0}{S_0-\bar D}\Bigr)^2t_1+\sigma^2(T-t_1).
\]
\end{proposition}
\begin{proof}
Avant $t_1$, l'action séquestrée $S^*=S-\bar D e^{rt}$ suit les
mouvements de l'action entière, $dS^*=dS=\sigma S\,dW$, de sorte que sa volatilité relative vaut $\sigma S/S^*$~; après
$t_1$, les deux coïncident. En additionnant les variances des deux périodes, avec $S/S^*$ figé
à sa valeur initiale, on obtient la formule.
\end{proof}

\begin{example}[Un dividende de 4]\label{ex:dv:dividends-borrow-and-forwards:adj}
Action à 100, options à un an, $r=3\%$, un dividende de 4 dans six mois, modèle spot
à volatilité 25~\%. Les modèles séquestré et proportionnel à 25~\% impliquent une volatilité implicite plate de 25.00~;
le modèle spot donne des prix qui impliquent 25.57 pour le strike 70 et 25.49 pour le 130~; la
volatilité effective de la proposition vaut 25.52
(voir la \cref{fig:dv:dividends-borrow-and-forwards:models}). Un demi-point de volatilité vaut
0.2 de prime à la monnaie~: un desk qui mélange les modèles d'un système à l'autre le comptabilise comme
un gain ou une perte.
\end{example}

\begin{omfigure}
\begin{tikzpicture}
\begin{axis}[width=10.5cm,height=4.8cm,
  xlabel={strike},ylabel={volatilité implicite (\%)},
  tick label style={font=\small},label style={font=\small},
  xmin=70,xmax=130,ymin=24.9,ymax=25.7,grid=major,grid style={gray!25},
  yticklabel style={/pgf/number format/.cd,fixed,precision=1},
  legend columns=2,legend style={font=\footnotesize,at={(0.5,-0.3)},anchor=north,draw=none,fill=none,/tikz/every even column/.append style={column sep=8pt}},
  table/col sep=comma]
\addplot[omDef,very thick] table[x=strike,y=escrowed]{figdata/derivatives/05-dividends-borrow-and-forwards/model_smiles.csv};
\addplot[omThm,very thick,mark=*,mark size=1.5pt] table[x=strike,y=spot]{figdata/derivatives/05-dividends-borrow-and-forwards/model_smiles.csv};
\addplot[omProp,thick,dashed] table[x=strike,y=adjusted]{figdata/derivatives/05-dividends-borrow-and-forwards/model_smiles.csv};
\legend{séquestré ou proportionnel à 25~\%,modèle spot à 25~\%,séquestré à la volatilité effective}
\end{axis}
\end{tikzpicture}

\omcaption{Options d'achat à un an sur une action à 100 avec un dividende de 4 dans six mois,
évalués par trois modèles de dividendes avec la même volatilité de 25~\% et relus en
volatilités implicites de Black sur le forward commun. Le modèle spot dépasse les autres d'un demi-point
et présente un léger skew~; la volatilité effective de
la \cref{prop:dv:dividends-borrow-and-forwards:adj} rejoint celui-ci à la monnaie. Données~:
le tutoriel.}\label{fig:dv:dividends-borrow-and-forwards:models}
\end{omfigure}

\begin{definition}[Modèle de dividendes mixte]\label{def:dv:dividends-borrow-and-forwards:mixed}
Un \emph{modèle de dividendes mixte}\index{modèle de dividendes mixte} traite les dividendes proches comme
des montants en numéraire, annoncés ou prévus avec confiance, et les dividendes au-delà d'un horizon
(un à deux ans) comme proportionnels au cours de l'action, en les mélangeant
sur une période de transition.
\end{definition}

Le mélange porte sur la dynamique. Avec des dividendes en numéraire, le forward évolue un pour un
avec le spot, $\partial F/\partial S=1/P(0,T)$~; avec des \omterm{def:dv:dividends-borrow-and-forwards:prop}{dividendes proportionnels}, il
évolue en proportion, $\partial F/\partial S=F/S$. Pour une option à cinq ans sur une action
versant un rendement de 4~\%, les deux deltas diffèrent de presque un cinquième, et il en va de même pour la couverture. Une société
dont l'action se réduit de moitié garde rarement le même dividende en numéraire, ce qui explique pourquoi les dividendes lointains
sont modélisés comme proportionnels~; les dividendes proches sont déjà déclarés.

\section{Déduire les forwards, les dividendes et l'emprunt de la parité}

Pour les options européennes, la parité call-put, $C-P=P(0,T)(F-K)$, vaut pour tout strike.
Elle est linéaire en fonction du strike, de pente $-P(0,T)$ et d'ordonnée à l'origine $P(0,T)F$, donc une
régression sur l'ensemble des strikes permet de lire ces deux grandeurs sur la chaîne, sans aucun modèle.

\begin{method}[Lire la courbe à terme sur la chaîne]\label{met:dv:dividends-borrow-and-forwards:read}
Pour chaque échéance~:
\begin{enumerate}
\item régresser $C_K-P_K$ sur $K$ pour les strikes proches de la monnaie~: la pente vaut
$-P(0,T)$, l'ordonnée à l'origine $P(0,T)F$~;
\item le carry implicite $S-FP(0,T)$ est la valeur actuelle des dividendes et du coût
de l'emprunt jusqu'à $T$~; ses accroissements entre échéances donnent la structure par terme des dividendes~;
\item si l'on dispose d'une prévision de dividendes, ce qui reste du carry est l'emprunt, la
définition suivante~; si l'emprunt est connu (collatéral général, GC), ce qui reste est la
prévision de dividendes du marché.
\end{enumerate}
Pour les options américaines sur action individuelle, la parité se réduit à un couple d'inégalités~; il faut
utiliser des options européennes sur indice, ou retrancher d'abord la prime d'exercice anticipé (chapitre~6).
\end{method}

\begin{definition}[Dividende implicite et taux d'emprunt implicite]\label{def:dv:dividends-borrow-and-forwards:implied}
Le \emph{dividende implicite}\index{dividende implicite} d'une échéance est la valeur actuelle
des dividendes jusqu'à cette échéance qui rend le forward déduit des options, avec
le coût d'emprunt connu, cohérent avec le spot. Le \emph{taux d'emprunt
implicite}\index{taux d'emprunt implicite} est le niveau des frais d'emprunt $\ell$ qui rend le forward
cohérent avec le spot et une prévision de dividendes~:
$\ell=-\ln\bigl(FP(0,T)/(S-\bar D)\bigr)/T$.
\end{definition}

\begin{example}[L'action difficile à emprunter]\label{ex:dv:dividends-borrow-and-forwards:htb}
La chaîne à un mois de l'introduction, cinq strikes de 45 à 55 cotés au centime,
régressée comme dans la \cref{met:dv:dividends-borrow-and-forwards:read} donne un forward de
48.74 et un facteur d'actualisation de 0.99680 (un taux de 3.90~\%, contre 4~\% réel~: les
centimes d'arrondi). Sans dividende à venir, le carry de 1.42 est entièrement de l'emprunt~: 35.0~\% par
an. Évalué à partir du forward naïf $50e^{0.04\times30/365}=50.16$ à la volatilité de 60~\%
de la chaîne, le call vaudrait 3.50 au lieu de 2.78, et le put 3.34 au lieu de
4.04. Le «~skew~» de 72 contre 47 tient au mauvais forward, et non à la vision que le marché
a de la volatilité.
\end{example}

Vendre le put, acheter le call et vendre à découvert l'action (le reversal du
livre~1, chapitre~25) semble rapporter 1.42 en un mois~; cela suppose d'emprunter l'action
pendant un mois à 35~\%, ce qui coûte 1.44. Les options sur actions difficiles à emprunter sont là où
la parité call-put est mise en défaut de la façon la plus visible lorsqu'on la mesure à partir du spot, et des études académiques concluent que
l'ampleur de la violation apparente suit les frais du marché du prêt de titres
(Ofek, Richardson et Whitelaw).

\begin{omfigure}
\begin{tikzpicture}
\begin{axis}[width=10.5cm,height=4.8cm,ybar,bar width=10pt,
  xlabel={échéance (années)},ylabel={carry par trimestre},
  tick label style={font=\small},label style={font=\small},
  xmin=0.1,xmax=2.15,ymin=0.5,ymax=0.9,ymajorgrids,grid style={gray!25},
  xtick={0.25,0.5,0.75,1,1.25,1.5,1.75,2},
  legend columns=2,legend style={font=\footnotesize,at={(0.5,-0.36)},anchor=north,draw=none,fill=none,/tikz/every even column/.append style={column sep=8pt}},
  table/col sep=comma]
\addplot[fill=omDef!55,draw=omDef] table[x=t,y=implied]{figdata/derivatives/05-dividends-borrow-and-forwards/carry_steps.csv};
\addplot[fill=omThm!45,draw=omThm] table[x=t,y=true]{figdata/derivatives/05-dividends-borrow-and-forwards/carry_steps.csv};
\legend{lu sur la parité (cotations $\pm$0.02),réel~: VA du dividende plus l'emprunt}
\end{axis}
\end{tikzpicture}

\omcaption{Valeur actuelle du carry (dividende plus coût de l'emprunt) entre des échéances
trimestrielles consécutives, retrouvée par la régression de parité à partir d'une chaîne synthétique dont
les prix milieu portent jusqu'à deux centimes de bruit sur neuf strikes. Le forward cumulé est
retrouvé à environ un centime près~; l'accroissement de chaque trimestre seulement à environ 0.08. Données~: le
tutoriel.}\label{fig:dv:dividends-borrow-and-forwards:carry}
\end{omfigure}

\section{Risque de dividende et swaps de dividendes}

\begin{definition}[Risque de dividende et swap de dividendes]\label{def:dv:dividends-borrow-and-forwards:risk}
Le \emph{risque de dividende}\index{risque de dividende} est la sensibilité d'une position à une variation
des dividendes attendus avant son échéance~: un forward, un call acheteur ou un put vendeur
perd de la valeur lorsque les dividendes attendus augmentent. Un \emph{swap de dividendes}\index{swap de dividendes} est un
contrat de gré à gré qui échange, à l'échéance, les dividendes effectivement versés sur
une action ou un indice pendant une période contre un montant fixe convenu à la conclusion~; la
version cotée est le contrat à terme sur dividendes du livre~1, chapitre~22.
\end{definition}

Le \omterm{def:dv:dividends-borrow-and-forwards:risk}{risque de dividende} est faible pour une option à un mois et dominant pour une option à cinq ans. Il
s'accumule au bilan des émetteurs de produits structurés~: une banque qui vend des notes
versant une fonction d'un indice se couvre en détenant l'indice, perçoit ses dividendes et
a intégré ces dividendes dans le prix des notes~; si les dividendes baissent, la couverture rapporte
moins que promis. Les praticiens décrivent ces émetteurs comme structurellement longs de dividendes et
comme vendeurs de \omterm{def:dv:dividends-borrow-and-forwards:risk}{swaps de dividendes} et de contrats à terme pour transférer le risque (chapitre~18). Les acheteurs
sont des investisseurs qui veulent une exposition aux dividendes sans exposition aux actions.

\begin{dated}{2026-09}{Contrats à terme sur dividendes d'indice et 2020}\label{dat:dv:dividends-borrow-and-forwards:divfut}
Les contrats à terme sur dividendes d'indice de référence de la zone euro, cotés sur Eurex, ont une valeur de contrat de
EUR~100 par point d'indice et un pas de cotation de 0.1 point (EUR~10), se liquident en espèces sur les
dividendes bruts cumulés versés par les membres de l'indice pendant l'année du contrat, et proposent
des échéances annuelles jusqu'à une dizaine d'années. Le 27 mars 2020, la Banque centrale européenne
a recommandé aux banques qu'elle supervise de ne verser aucun dividende au moins jusqu'au 1 octobre
2020, prolongé ensuite jusqu'en janvier 2021. Un compte rendu de praticien d'avril 2020 faisait état de
dividendes d'indice 2020 négociés à 55~\% de ce que les modèles avaient prévu, et de dividendes 2021 à
70~\%, les émetteurs vendant des dividendes dans un marché en baisse.
\end{dated}

\begin{omfigure}
\begin{tikzpicture}[font=\small,
  bx/.style={draw=omDef,thick,rounded corners=2pt,fill=omDef!6,align=center,text width=2.6cm,inner sep=4pt}]
\node[bx] (inv) at (0,0) {investisseurs dans les\\notes structurées};
\node[bx] (iss) at (5.6,0) {émetteur~: détient\\l'indice en couverture,\\long dividendes};
\node[bx] (buy) at (11.2,0) {acheteurs de dividendes~:\\fonds, fonds\\de pension};
\node[bx] (co) at (5.6,-2.6) {membres de l'indice\\versent les dividendes};
\draw[-{Stealth},thick] (inv) -- (iss);
\draw[-{Stealth},thick] (iss) -- (buy);
\node[font=\footnotesize,align=center,anchor=south] at (2.8,0.1) {achat de notes\\intégrant les\\dividendes};
\node[font=\footnotesize,align=center,anchor=south] at (8.4,0.1) {vente de swaps\\et de futures\\sur dividendes};
\draw[-{Stealth},thick,omThm] (co) -- node[right,font=\footnotesize]{dividendes} (iss);
\end{tikzpicture}

\omcaption{Où va le \omterm{def:dv:dividends-borrow-and-forwards:risk}{risque de dividende}. L'émetteur de notes indexées sur actions se couvre avec le
sous-jacent, perçoit ses dividendes et a promis leur valeur espérée aux porteurs de
notes~; il transfère le risque en vendant à terme des dividendes aux investisseurs qui
les veulent.}\label{fig:dv:dividends-borrow-and-forwards:flows}
\end{omfigure}

\section{Tutoriel~: une courbe à terme à partir d'une chaîne}\label{tut:dv:dividends-borrow-and-forwards:tut}

\begin{tutorial}
\textbf{Objectif.} Lire les forwards, les facteurs d'actualisation, le carry et l'emprunt à partir d'une chaîne
synthétique, et comparer les trois modèles de dividendes sur une même option. \textbf{État final~:}
\cref{fig:dv:dividends-borrow-and-forwards:models,fig:dv:dividends-borrow-and-forwards:carry}
et les résultats de l'\cref{ex:dv:dividends-borrow-and-forwards:htb}.
\begin{tutsteps}
\item \textbf{Parité par régression}, puis carry, emprunt et strip de dividendes~:
\omcode{code/firm/divfwd/firm_divfwd.py}{34}{58}{Forward et facteur d'actualisation à partir d'une chaîne~; carry implicite, emprunt et strip de dividendes.}\label{lst:dv:dividends-borrow-and-forwards:parity}
\item \textbf{Le modèle spot} par quadrature de Gauss--Hermite sur l'action juste avant la
date de détachement, et la volatilité effective de l'action séquestrée~:
\omcode{code/derivatives/05-dividends-borrow-and-forwards/python/dv_dividends.py}{35}{53}{Le modèle spot, et la volatilité effective du modèle à dividendes séquestrés.}\label{lst:dv:dividends-borrow-and-forwards:spot}
\item \textbf{Exécuter} \texttt{dv\_dividends.curve\_from\_chain()},
\texttt{htb\_results()} et \texttt{fig\_dividends.py}.
\end{tutsteps}
\textbf{Pour aller plus loin.} Régresser sur les cinq strikes les plus proches du forward uniquement, et
sur tous les strikes pondérés par l'inverse de la largeur de cotation, et comparer le bruit du carry
trimestriel~; puis, donner à l'action deux dividendes avant l'échéance et étendre la volatilité
effective aux deux.
\end{tutorial}

\section{Construction~: la courbe à terme sur actions}\label{bld:dv:dividends-borrow-and-forwards:divfwd}

\begin{build}
\bfield{Rôle} Chaque option sur action de la firme miniature est évaluée sur un forward qui
provient de cette courbe~; dans la bibliothèque de tarification du chapitre~28, c'est ce que
calcule \texttt{MarketData.forward}. \bfield{Interface}
\texttt{ForwardCurve(spot, rate, cash, proportional,
borrow).forward(t)}~; \texttt{implied\_forward(strikes, calls, puts) -> (F, df)}~;
\texttt{implied\_carry}~; \texttt{implied\_borrow(spot, fwd, df, t, pv\_dividends)}~;
\texttt{strip\_dividends(spot, expiries, forwards, dfs)}.
\bfield{Règles} Les dividendes strictement postérieurs à aujourd'hui et antérieurs ou égaux à la date de livraison comptent~;
les \omterm{def:dv:dividends-borrow-and-forwards:prop}{dividendes proportionnels} multiplient le forward~; la régression utilise les cotations européennes (ou
leurs équivalents désaméricanisés).
\bfield{Tests d'acceptation} \texttt{code/firm/divfwd/tests/}~: le forward se réplique par
l'argument du cash-and-carry~; la régression retrouve un forward et un facteur d'actualisation connus
exactement à partir de cotations sans bruit~; l'emprunt implicite d'une chaîne propre égale les frais utilisés
pour la construire~; les strips somment au carry total.
\bfield{Pour aller plus loin} Ajuster conjointement une courbe d'emprunt lisse et une courbe de dividendes sur l'ensemble des échéances,
avec une pénalité sur la rugosité, et signaler les échéances dont les résidus de parité sont aberrants.
\end{build}

\begin{omsources}
\item E.~G.~Haug, J.~Haug et A.~Lewis, ``Back to basics: a new approach to the
discrete dividend problem'', \emph{Wilmott} (septembre 2003) 37--47.
\item E.~Ofek, M.~Richardson et R.~F.~Whitelaw, ``Limited arbitrage and short sales
restrictions: evidence from the options markets'', \emph{Journal of Financial
Economics} 74 (2004) 305--342.
\item Banque centrale européenne, recommandation ECB/2020/19 du 27 mars 2020 sur les distributions
de dividendes pendant la pandémie de COVID-19.
\item Eurex, EURO STOXX~50 Index Dividend Futures, spécifications des contrats.
\item E.~Barthe, ``Hedging structured products during the corona crisis: where did it
go wrong?'', Structured Retail Products, 30 avril 2020.
\end{omsources}

\section{Exercices}

\begin{exercise}[$\star$]\label{exo:dv:dividends-borrow-and-forwards:1}
Une action à 100 verse des dividendes de 1 dans trois mois et dans neuf mois~; $r=3\%$, sans frais
d'emprunt. Donner le terme à un an.
\end{exercise}

\begin{exercise}[$\star$]\label{exo:dv:dividends-borrow-and-forwards:2}
À une échéance donnée, $C-P$ vaut 7.02 au strike 95 et $-2.83$ au strike 105. Donner le
facteur d'actualisation et le terme.
\end{exercise}

\begin{exercise}[$\star$]\label{exo:dv:dividends-borrow-and-forwards:3}
Vérifier la volatilité effective de 25.52~\% de l'\cref{ex:dv:dividends-borrow-and-forwards:adj}.
\end{exercise}

\begin{exercise}[$\star\star$]\label{exo:dv:dividends-borrow-and-forwards:4}
Pourquoi la parité call-put ne donne-t-elle que des bornes pour les options américaines sur actions individuelles, et dans
quel sens se trompe-t-on lorsque la relation européenne est utilisée à proximité d'une date de détachement~?
\end{exercise}

\begin{exercise}[$\star\star$]\label{exo:dv:dividends-borrow-and-forwards:5}
Une action à 100 verse un dividende valant 4~\% de son prix dans six mois~; $r=3\%$, horizon d'un an.
Donner $\partial F/\partial S$ lorsque le dividende est modélisé comme un montant en espèces égal à 4, et
lorsqu'il est modélisé comme proportionnel.
\end{exercise}

\begin{exercise}[$\star\star$]\label{exo:dv:dividends-borrow-and-forwards:6}
Un call à la monnaie de trois ans sur l'action de la \cref{fig:dv:dividends-borrow-and-forwards:curve}
(dividendes trimestriels de 0.60, $r=3\%$, emprunt à 0.5~\%, volatilité de 20~\%) est réévalué après une
baisse de 10~\% de chaque dividende futur. De combien sa valeur change-t-elle~?
\end{exercise}

\begin{exercise}[$\star\star\star$]\label{exo:dv:dividends-borrow-and-forwards:7}
\emph{Codage.} À partir de \texttt{synthetic\_chain()}, calculer le terme implicite et
le facteur d'actualisation de l'échéance à un an, ainsi que l'emprunt implicite compte tenu des dividendes connus,
puis comparer avec les 0.5~\% utilisés pour construire la chaîne.
\end{exercise}

\begin{exercise}[$\star\star\star$]\label{exo:dv:dividends-borrow-and-forwards:8}
\emph{Trouver l'erreur.} «~Sur cette action difficile à emprunter, les puts à un mois de strike 50 se traitent à 72
de volatilité et les calls à 47. En vendant des puts et en achetant des calls, couverts en delta, on capte 25
points de volatilité.~»
\end{exercise}

\section{Problème~: L'action difficile à emprunter}

\begin{problem}[{Problème du week-end --- lire l'emprunt dans les options}]\label{pb:dv:dividends-borrow-and-forwards:1}
L'action à l'ouverture se traite à 50 et ne verse aucun dividende le mois suivant~; le
taux de financement est de 4~\%. Ses options européennes à un mois sont cotées (prix milieu)~: call 45 à 5.41,
put à 1.68~; call 47.5 à 3.94, put à 2.71~; call 50 à 2.78, put à 4.04~; call 52.5 à 1.90, put à 5.65~; call 55
à 1.26, put à 7.50. La volatilité à la monnaie de la chaîne est de 60~\%.

\medskip\noindent\textbf{Partie I --- Le forward.}
\begin{enumerate}
\item Calculer $C-P$ à chaque strike.
\item Régresser ce montant sur le strike~: donner le facteur d'actualisation et le terme.
\item Donner le taux implicite, et expliquer l'écart avec 4~\%.
\item Donner le portage implicite.
\item Donner le \omterm{def:dv:dividends-borrow-and-forwards:implied}{taux d'emprunt implicite}.
\end{enumerate}

\medskip\noindent\textbf{Partie II --- Le mauvais forward.}
\begin{enumerate}[resume]
\item Donner le terme naïf à partir du spot et du taux.
\item Pricer le call et le put de strike 50 à 60~\% sur le terme naïf.
\item De combien sont-ils mal valorisés~?
\item Donner leurs volatilités implicites calculées sur le terme naïf.
\item Pourquoi les deux volatilités implicites diffèrent-elles, et quel terme les rend-il égales~?
\end{enumerate}

\medskip\noindent\textbf{Partie III --- L'arbitrage apparent.}
\begin{enumerate}[resume]
\item Que rapporte apparemment le renversement (vente à découvert de l'action, vente de put, achat de call)~?
\item Combien coûte l'emprunt de l'action pour le mois~?
\item Y a-t-il un arbitrage une fois l'emprunt pris en compte~?
\item Que peut faire un détenteur de l'action avec les options plutôt que de la prêter, et
que lui rapporte cela~?
\item Comment un desk incapable d'emprunter l'action exprimerait-il une vue selon laquelle elle baissera~?
\end{enumerate}

\medskip\noindent\textbf{Partie IV --- Jugement.}
\begin{enumerate}[resume]
\item Que se passe-t-il pour les options si les frais d'emprunt tombent à 5~\% du jour au lendemain~?
\item Pourquoi les options cotées sur de telles actions sont-elles américaines, et qu'est-ce que cela change~?
\item Quel terme un système d'options doit-il utiliser pour cette action~?
\item Énoncer le \emph{résultat nommé}~: l'emprunt implicite et l'erreur de pricing naïve sur
le call de strike 50.
\item En une phrase~: que mesure la parité call-put sur une action difficile à emprunter~?
\end{enumerate}
\end{problem}

\section{Questions d'entretien}

\begin{interviewq}[$\star$ \iqroles{trader, chercheur}]\label{iq:dv:dividends-borrow-and-forwards:1}
Comment obtenez-vous le terme d'une action à partir de ses prix d'options, sans modèle~?
\end{interviewq}

\begin{interviewq}[$\star$ \iqroles{trader}]\label{iq:dv:dividends-borrow-and-forwards:2}
Le put à la monnaie d'une action se traite nettement au-dessus du call au même strike. Donnez deux
raisons.
\end{interviewq}

\begin{interviewq}[$\star\star$ \iqroles{chercheur, banque}]\label{iq:dv:dividends-borrow-and-forwards:3}
Dividendes en espèces ou proportionnels~: lequel utilisez-vous pour une option à un mois et pour une
option à cinq ans, et qu'est-ce qui change dans la couverture~?
\end{interviewq}

\begin{interviewq}[$\star\star$ \iqroles{trader, risque}]\label{iq:dv:dividends-borrow-and-forwards:4}
Qu'est-ce que le \omterm{def:dv:dividends-borrow-and-forwards:risk}{risque de dividende}, qui est naturellement long de ce risque, et que lui est-il arrivé en 2020~?
\end{interviewq}

\begin{interviewq}[$\star\star$ \iqroles{développeur, chercheur}]\label{iq:dv:dividends-borrow-and-forwards:5}
Comment intégrez-vous un dividende en espèces discret dans un pricer à arbres pour options américaines, et
qu'est-ce qui ne va pas avec les approches les plus évidentes~?
\end{interviewq}

\begin{interviewq}[$\star\star\star$ \iqroles{chercheur}]\label{iq:dv:dividends-borrow-and-forwards:6}
La même volatilité de 25~\% est donnée au \omterm{def:dv:dividends-borrow-and-forwards:escrowed}{modèle à dividendes séquestrés} dans un système et au modèle spot
dans un autre. Lequel donne le prix de call le plus élevé, de combien environ, et pourquoi~?
\end{interviewq}
